% 03-phan-tich-thiet-ke.tex — Chương 3: phân tích và thiết kế hệ thống.
% Mỗi nhóm chức năng một mục (3.2 Tài khoản, các nhóm sau thêm 3.3, 3.4...).
% Nguồn: docs/design/use-cases.md, docs/specs/<nhóm>.md, docs/diagrams/*.puml (ảnh: images/diagrams/).
\chapter{Phân tích và thiết kế hệ thống}

\section{Tổng quan về các ca sử dụng}

Hệ thống có bảy tác nhân. \textit{Người dùng} là tác nhân chung, đại diện cho mọi người đã hoặc
sắp có tài khoản; sáu tác nhân còn lại kế thừa từ tác nhân này: Người tham dự, Ban tổ chức, Diễn giả,
Nhà cung cấp, Nhân viên vận hành và Quản trị viên. Danh sách ca sử dụng được chốt gồm 57 ca, đánh
mã theo thứ tự ưu tiên: các ca có mã nhỏ nằm ở luồng lõi (từ tạo hội nghị đến công bố lịch) và được
xây dựng trước.

Khóa luận trình bày phần phân tích và thiết kế theo từng nhóm chức năng, đúng thứ tự xây dựng hệ
thống. Mỗi nhóm gồm: phạm vi và tác nhân, quy tắc nghiệp vụ, đặc tả ca sử dụng, các biểu đồ hoạt
động, tuần tự, trạng thái, thiết kế dữ liệu và thiết kế API.

\chuaviet[Biểu đồ ca sử dụng tổng quát và danh sách 57 ca sử dụng sẽ được bổ sung.]

% =====================================================================================
\section{Nhóm chức năng Tài khoản}
\label{sec:tai-khoan}

\subsection{Phạm vi và tác nhân}

Nhóm Tài khoản là điểm vào của mọi tác nhân: trước khi dùng bất kỳ chức năng nào, người dùng phải có
tài khoản và đăng nhập. Nhóm gồm hai ca sử dụng của tác nhân Người dùng, được thể hiện trong
Hình~\ref{fig:uc-nguoi-dung}.

\begin{figure}[H]
  \centering
  \includegraphics[width=0.55\textwidth]{diagrams/uc-nguoi-dung.png}
  \caption{Biểu đồ ca sử dụng của tác nhân Người dùng}
  \label{fig:uc-nguoi-dung}
\end{figure}

Hình~\ref{fig:uc-nguoi-dung} cho thấy sáu tác nhân cụ thể đều kế thừa tác nhân Người dùng, nên mọi
tác nhân đều dùng chung cơ chế đăng ký và đăng nhập. Ca UC41 (quản lý hồ sơ cá nhân) thuộc nhóm hỗ
trợ, được xây dựng ở giai đoạn sau. Hai ca UC01 và UC02 được chia thành các luồng con như
Bảng~\ref{tab:tk-pham-vi}.

\begin{table}[H]
  \centering
  \caption{Các ca sử dụng của nhóm Tài khoản}
  \label{tab:tk-pham-vi}
  \begin{tabularx}{\textwidth}{lp{4cm}L}
    \toprule
    \textbf{Mã} & \textbf{Tên} & \textbf{Luồng con} \\
    \midrule
    UC01 & Đăng ký tài khoản & Đăng ký; xác thực email; gửi lại email xác thực \\
    UC02 & Đăng nhập / đăng xuất & Đăng nhập; làm mới phiên; đăng xuất; quên và đặt lại mật khẩu \\
    \bottomrule
  \end{tabularx}
\end{table}

``Làm mới phiên'' không phải thao tác của người dùng: giao diện tự thực hiện khi phiên làm việc hết
hạn hoặc khi người dùng mở lại trang. Việc quản trị viên duyệt, khóa tài khoản (UC20) và quản lý vai
trò, quyền thuộc các nhóm chức năng sau.

Hệ thống có năm vai trò. Ba vai trò được tự đăng ký; hai vai trò còn lại do quản trị viên tạo
(Bảng~\ref{tab:tk-vai-tro}). Diễn giả không phải một vai trò riêng: một người dùng bất kỳ trở thành
diễn giả khi được Ban tổ chức mời trình bày ở một phiên.

\begin{table}[H]
  \centering
  \caption{Vai trò và trạng thái sau khi xác thực email}
  \label{tab:tk-vai-tro}
  \begin{tabularx}{\textwidth}{lLcl}
    \toprule
    \textbf{Mã vai trò} & \textbf{Tên} & \textbf{Tự đăng ký} & \textbf{Sau xác thực email} \\
    \midrule
    \code{Attendee}  & Người tham dự      & Có    & Hoạt động (\code{Active}) \\
    \code{Organizer} & Ban tổ chức        & Có    & Chờ duyệt (\code{PendingApproval}) \\
    \code{Supplier}  & Nhà cung cấp       & Có    & Chờ duyệt (\code{PendingApproval}) \\
    \code{Staff}     & Nhân viên vận hành & Không & -- \\
    \code{Admin}     & Quản trị viên      & Không & -- \\
    \bottomrule
  \end{tabularx}
\end{table}

Ban tổ chức và Nhà cung cấp phải được quản trị viên duyệt vì hai vai trò này tạo dữ liệu ảnh hưởng
tới người khác (hội nghị, đơn dịch vụ); Người tham dự được hoạt động ngay sau khi xác thực email.

\subsection{Quy tắc nghiệp vụ}

Các quy tắc của nhóm Tài khoản được đánh mã BR01--BR15 (Bảng~\ref{tab:tk-quy-tac}) và được kiểm tra
ở cả giao diện lẫn máy chủ; máy chủ luôn kiểm tra lại, không tin dữ liệu gửi từ trình duyệt.

\begin{longtable}{p{1.3cm}p{3.2cm}p{9.9cm}}
  \caption{Quy tắc nghiệp vụ của nhóm Tài khoản}\label{tab:tk-quy-tac}\\
  \toprule
  \textbf{Mã} & \textbf{Nội dung} & \textbf{Quy tắc} \\
  \midrule
  \endfirsthead
  \toprule
  \textbf{Mã} & \textbf{Nội dung} & \textbf{Quy tắc} \\
  \midrule
  \endhead
  BR01 & Email & Bắt buộc, đúng định dạng, tối đa 256 ký tự; cắt khoảng trắng và đổi về chữ thường
  trước khi lưu, so sánh; duy nhất trong hệ thống. \\
  BR02 & Mật khẩu & 8--128 ký tự, có ít nhất một chữ cái và một chữ số. Chỉ lưu bản băm. \\
  BR03 & Họ tên & Bắt buộc, 2--100 ký tự sau khi cắt khoảng trắng. \\
  BR04 & Đơn vị & Tối đa 200 ký tự; bắt buộc với Ban tổ chức và Nhà cung cấp. \\
  BR05 & Vai trò khi đăng ký & Chỉ nhận Người tham dự, Ban tổ chức, Nhà cung cấp. \\
  BR06 & Token một lần & Token trong liên kết email là chuỗi ngẫu nhiên 32 byte; CSDL chỉ lưu
  bản băm SHA-256. Token đã dùng, hết hạn hoặc đã có token mới hơn thì không còn hiệu lực. \\
  BR07 & Token không đi qua hàng đợi & Sự kiện chỉ mang mã người dùng; token được sinh ngay lúc
  gửi thư, chuỗi gốc chỉ nằm trong email. \\
  BR08 & Gửi lại email & Cách lần gửi trước cùng loại ít nhất 60 giây; giao diện khóa nút trong
  thời gian này. \\
  BR09 & Không lộ email & Chức năng gửi lại email xác thực và quên mật khẩu luôn trả cùng một thông
  báo, dù email có tồn tại hay không. \\
  BR10 & Đăng nhập & Chỉ tài khoản đang hoạt động được cấp phiên. Sai email hoặc sai mật khẩu trả
  cùng một thông báo; trạng thái chỉ được kiểm tra sau khi mật khẩu đúng. \\
  BR11 & Token phiên & Access token (JWT) hiệu lực 15 phút, giao diện giữ trong bộ nhớ. Refresh token
  hiệu lực 7 ngày, nằm trong cookie \code{HttpOnly}. \\
  BR12 & Xoay vòng refresh token & Mỗi lần làm mới, token cũ bị đánh dấu đã dùng và cấp token mới.
  Token đã dùng bị gửi lại thì thu hồi mọi phiên của người đó. \\
  BR13 & Đăng xuất & Thu hồi refresh token hiện tại và xóa cookie. \\
  BR14 & Đặt lại mật khẩu & Thành công thì thu hồi mọi refresh token (đăng xuất khỏi mọi thiết bị). \\
  BR15 & Giới hạn tần suất & Theo địa chỉ IP: đăng nhập, đăng ký 5 lần/phút; gửi lại email xác thực,
  quên mật khẩu 3 lần/phút. Không khóa tài khoản khi nhập sai nhiều lần. \\
  \bottomrule
\end{longtable}

Hai quy tắc BR09 và BR10 nhằm chống dò tài khoản: kẻ tấn công không thể dựa vào thông báo trả về để
biết một email đã đăng ký hay tài khoản đang ở trạng thái nào. Riêng chức năng đăng ký buộc phải báo
``email đã được sử dụng'' để người dùng biết cách xử lý; rủi ro này được giảm bằng giới hạn tần suất
BR15.

\subsection{Đặc tả ca sử dụng UC01 -- Đăng ký tài khoản}

\begin{table}[H]
  \centering
  \caption{Thông tin chung của UC01}
  \label{tab:uc01}
  \begin{tabularx}{\textwidth}{p{3.5cm}L}
    \toprule
    Tác nhân & Người dùng (chưa có tài khoản) \\
    Mục tiêu & Tạo tài khoản Người tham dự, Ban tổ chức hoặc Nhà cung cấp \\
    Tiền điều kiện & Chưa đăng nhập \\
    Hậu điều kiện & Tài khoản ở trạng thái chưa xác thực; email xác thực đã được gửi \\
    Kích hoạt & Bấm ``Đăng ký'' trên thanh tiêu đề hoặc trang đăng nhập \\
    \bottomrule
  \end{tabularx}
\end{table}

\textbf{Luồng chính -- Đăng ký:}
\begin{enumerate}
  \item Người dùng mở trang đăng ký, chọn vai trò (mặc định Người tham dự). Nếu chọn Ban tổ chức hoặc
  Nhà cung cấp, giao diện hiện thêm ô ``Đơn vị'' bắt buộc và ghi chú cần quản trị viên duyệt.
  \item Người dùng nhập họ tên, email, mật khẩu, nhập lại mật khẩu (và đơn vị), bấm ``Đăng ký''.
  \item Giao diện kiểm tra BR01--BR04; lỗi hiện ngay dưới ô nhập.
  \item Hệ thống kiểm tra lại BR01--BR05 và kiểm tra email chưa tồn tại.
  \item Hệ thống băm mật khẩu, tạo tài khoản chưa xác thực và ghi sự kiện ``gửi email xác thực'' trong
  cùng một giao dịch.
  \item Giao diện chuyển sang trang ``Kiểm tra email'', hiển thị địa chỉ vừa đăng ký.
  \item Phía sau, hệ thống sinh token xác thực, lưu bản băm và gửi email chứa liên kết xác thực.
\end{enumerate}

\textbf{Luồng con -- Xác thực email:} người dùng bấm liên kết trong email; giao diện gửi token lên hệ
thống. Hệ thống tìm token theo bản băm, kiểm tra đúng mục đích, chưa dùng, chưa hết hạn (BR06), đánh
dấu token đã dùng rồi đổi trạng thái tài khoản theo vai trò (Bảng~\ref{tab:tk-vai-tro}).

\textbf{Luồng con -- Gửi lại email xác thực:} ở trang ``Kiểm tra email'' hoặc khi đăng nhập bị báo chưa
xác thực, người dùng bấm ``Gửi lại email''. Hệ thống áp dụng BR08 và làm token cũ hết hiệu lực khi
tạo token mới; phản hồi luôn là ``Nếu email chưa xác thực, chúng tôi đã gửi lại liên kết'' (BR09).

\textbf{Ngoại lệ:}
\begin{itemize}
  \item E1 -- Dữ liệu sai quy tắc: báo lỗi theo từng ô nhập.
  \item E2 -- Email đã tồn tại: báo ``Email đã được sử dụng'', gợi ý đăng nhập hoặc lấy lại mật khẩu.
  \item E3 -- Liên kết xác thực sai, hết hạn hoặc đã dùng: báo ``Liên kết không còn hiệu lực'' kèm ô
  nhập email để nhận liên kết mới.
  \item E4 -- Máy chủ thư lỗi: hệ thống tự gửi lại theo chính sách thử lại; việc đăng ký vẫn thành
  công và người dùng có thể bấm ``Gửi lại''.
  \item E5 -- Thao tác quá nhanh (vượt BR15): báo ``Thao tác quá nhanh, thử lại sau''.
\end{itemize}

\begin{figure}[H]
  \centering
  \includegraphics[height=0.72\textheight,keepaspectratio]{diagrams/g1-hd-dang-ky.png}
  \caption{Biểu đồ hoạt động đăng ký, xác thực và gửi lại email}
  \label{fig:hd-dang-ky}
\end{figure}

Hình~\ref{fig:hd-dang-ky} thể hiện ba luồng của UC01 trên hai làn Người dùng và Hệ thống. Vòng lặp
đầu tiên là vòng nhập liệu: hệ thống chỉ tạo tài khoản khi dữ liệu hợp lệ và email chưa tồn tại. Vòng
lặp thứ hai là vòng chờ email: mỗi lần người dùng bấm gửi lại, hệ thống chỉ gửi khi đã quá 60 giây
nhưng luôn trả cùng một thông báo. Cuối cùng, nhánh rẽ theo vai trò quyết định tài khoản chuyển sang
hoạt động ngay hay phải chờ quản trị viên duyệt.

\subsection{Đặc tả ca sử dụng UC02 -- Đăng nhập / đăng xuất}

\begin{table}[H]
  \centering
  \caption{Thông tin chung của UC02}
  \label{tab:uc02}
  \begin{tabularx}{\textwidth}{p{3.5cm}L}
    \toprule
    Tác nhân & Người dùng (cả sáu tác nhân kế thừa) \\
    Mục tiêu & Vào hệ thống đúng quyền của vai trò; thoát an toàn \\
    Tiền điều kiện & Đã có tài khoản \\
    Hậu điều kiện & Đăng nhập: có phiên làm việc; đăng xuất: phiên bị thu hồi \\
    Kích hoạt & Bấm ``Đăng nhập'', ``Đăng xuất'' hoặc mở trang cần đăng nhập \\
    \bottomrule
  \end{tabularx}
\end{table}

\textbf{Luồng chính -- Đăng nhập:}
\begin{enumerate}
  \item Người dùng mở trang đăng nhập (hoặc bị chuyển tới, kèm địa chỉ trang muốn vào).
  \item Người dùng nhập email, mật khẩu và bấm ``Đăng nhập''.
  \item Hệ thống tìm tài khoản theo email, so mật khẩu với bản băm (BR10).
  \item Hệ thống kiểm tra tài khoản đang hoạt động.
  \item Hệ thống tạo access token kèm danh sách quyền của vai trò và refresh token; lưu bản băm
  refresh token; trả access token trong nội dung phản hồi và refresh token trong cookie (BR11).
  \item Giao diện giữ access token trong bộ nhớ và chuyển tới trang người dùng muốn vào, hoặc trang
  đầu của vai trò.
\end{enumerate}

\textbf{Luồng con -- Làm mới phiên:} khi mở lại trang hoặc khi một yêu cầu bị từ chối vì access token
hết hạn, giao diện gọi chức năng làm mới; trình duyệt tự gửi cookie. Hệ thống kiểm tra refresh token
còn hiệu lực và tài khoản còn hoạt động, đánh dấu token cũ đã dùng, cấp cặp token mới (BR12); giao
diện gửi lại yêu cầu vừa bị từ chối. Làm mới thất bại thì giao diện xóa phiên và đưa về trang đăng nhập.

\textbf{Luồng con -- Đăng xuất:} hệ thống thu hồi refresh token trong cookie và xóa cookie (BR13);
giao diện xóa access token, xóa dữ liệu tạm và về trang chủ.

\textbf{Luồng con -- Quên mật khẩu:} người dùng nhập email. Tài khoản hoạt động hoặc chờ duyệt nhận
thư đặt lại mật khẩu (hiệu lực 30 phút); tài khoản chưa xác thực nhận lại thư xác thực; tài khoản bị
khóa hoặc không tồn tại không nhận gì. Phản hồi luôn là ``Nếu email tồn tại, chúng tôi đã gửi hướng
dẫn'' (BR09). Đặt mật khẩu mới thành công thì mọi phiên bị thu hồi (BR14).

\textbf{Ngoại lệ:}
\begin{itemize}
  \item E1 -- Sai email hoặc mật khẩu: ``Email hoặc mật khẩu không đúng''.
  \item E2 -- Chưa xác thực email: thông báo kèm nút ``Gửi lại email xác thực''.
  \item E3 -- Đang chờ duyệt: ``Tài khoản đang chờ quản trị viên duyệt''.
  \item E4 -- Bị khóa: ``Tài khoản đã bị khóa, liên hệ quản trị viên''.
  \item E5 -- Refresh token đã dùng bị gửi lại: thu hồi mọi phiên của tài khoản (BR12).
  \item E6 -- Liên kết đặt lại mật khẩu không hợp lệ: ``Liên kết không còn hiệu lực''.
  \item E7 -- Vượt giới hạn tần suất: ``Thao tác quá nhanh, thử lại sau''.
\end{itemize}

\begin{figure}[H]
  \centering
  \includegraphics[width=0.95\textwidth]{diagrams/g1-hd-dang-nhap.png}
  \caption{Biểu đồ hoạt động đăng nhập}
  \label{fig:hd-dang-nhap}
\end{figure}

Hình~\ref{fig:hd-dang-nhap} cho thấy thứ tự kiểm tra khi đăng nhập: hệ thống so mật khẩu trước, chỉ
khi mật khẩu đúng mới xét trạng thái tài khoản. Nhờ vậy người không biết mật khẩu không dò được tài
khoản đang chờ duyệt hay bị khóa. Ba trạng thái chưa xác thực, chờ duyệt, bị khóa đều dừng ở thông
báo riêng; chỉ tài khoản hoạt động mới được cấp token.

\begin{figure}[H]
  \centering
  \includegraphics[height=0.62\textheight,keepaspectratio]{diagrams/g1-hd-quen-mat-khau.png}
  \caption{Biểu đồ hoạt động quên và đặt lại mật khẩu}
  \label{fig:hd-quen-mat-khau}
\end{figure}

Trong Hình~\ref{fig:hd-quen-mat-khau}, bước phân nhánh theo trạng thái tài khoản xảy ra hoàn toàn
phía máy chủ; người dùng luôn nhận cùng một thông báo. Ở bước đặt mật khẩu mới, hệ thống kiểm tra
token trước khi lưu mật khẩu, và sau khi đổi thành công thì thu hồi mọi refresh token để phiên đăng
nhập trên các thiết bị khác (có thể của kẻ đã chiếm tài khoản) không còn dùng được.

\subsection{Trạng thái tài khoản}

\begin{figure}[H]
  \centering
  \includegraphics[width=0.6\textwidth]{diagrams/g1-trang-thai-tai-khoan.png}
  \caption{Biểu đồ trạng thái của tài khoản}
  \label{fig:trang-thai-tai-khoan}
\end{figure}

Hình~\ref{fig:trang-thai-tai-khoan} mô tả vòng đời của một tài khoản qua bốn trạng thái. Tài khoản
tự đăng ký bắt đầu ở trạng thái chưa xác thực; tài khoản do quản trị viên tạo (Nhân viên vận hành,
Quản trị viên) bắt đầu ngay ở trạng thái hoạt động. Các chuyển trạng thái duyệt, từ chối, khóa và mở
khóa do quản trị viên thực hiện ở UC20. Chỉ trạng thái hoạt động được phép đăng nhập.

\subsection{Biểu đồ tuần tự}

\begin{figure}[H]
  \centering
  \includegraphics[width=\textwidth]{diagrams/g1-td-dang-ky.png}
  \caption{Biểu đồ tuần tự đăng ký, gửi email và xác thực email}
  \label{fig:td-dang-ky}
\end{figure}

Hình~\ref{fig:td-dang-ky} chia quá trình đăng ký thành ba đoạn. Đoạn ``Đăng ký'' kết thúc ở bước 10:
máy chủ trả kết quả ngay sau khi lưu tài khoản và sự kiện vào cơ sở dữ liệu, chưa chờ gửi thư. Đoạn
``Gửi email'' chạy nền: sự kiện được chuyển qua hàng đợi RabbitMQ tới bộ xử lý gửi thư; bộ xử lý này
mới sinh token, lưu bản băm và gửi thư. Đoạn ``Xác thực email'' bắt đầu khi người dùng bấm liên kết;
hệ thống chỉ cần so bản băm của token nhận được với bản băm đã lưu. Cách tách này giúp việc đăng ký
không bị ảnh hưởng khi máy chủ thư chậm hoặc tạm ngừng.

\begin{figure}[H]
  \centering
  \includegraphics[height=0.75\textheight,keepaspectratio]{diagrams/g1-td-dang-nhap.png}
  \caption{Biểu đồ tuần tự đăng nhập, làm mới phiên và đăng xuất}
  \label{fig:td-dang-nhap}
\end{figure}

Hình~\ref{fig:td-dang-nhap} gồm ba đoạn. Khi đăng nhập, access token được trả trong nội dung phản
hồi còn refresh token nằm trong cookie \code{HttpOnly}, mã JavaScript không đọc được. Khi access token
hết hạn, giao diện tự gọi làm mới và gửi lại yêu cầu vừa bị từ chối, người dùng không phải đăng nhập
lại. Nhánh ``token đã dùng bị gửi lại'' là dấu hiệu refresh token bị đánh cắp: hệ thống thu hồi toàn bộ
phiên của người dùng đó. Khi đăng xuất, refresh token bị đánh dấu đã dùng và cookie bị xóa.

\subsection{Thiết kế dữ liệu}

Nhóm Tài khoản dùng ba bảng (Hình~\ref{fig:erd-tai-khoan}). Mọi bảng có khóa chính kiểu
\code{uniqueidentifier} và cột thời điểm tạo (không vẽ trong hình).

\begin{figure}[H]
  \centering
  \includegraphics[width=0.42\textwidth]{diagrams/g1-erd-tai-khoan.png}
  \caption{Lược đồ dữ liệu nhóm Tài khoản}
  \label{fig:erd-tai-khoan}
\end{figure}

Theo Hình~\ref{fig:erd-tai-khoan}, mỗi vai trò có nhiều người dùng, mỗi người dùng có nhiều token.
Bảng \code{Role} lưu danh sách quyền của vai trò dưới dạng JSON; quyền này được đưa vào access token
khi đăng nhập. Cột \code{Bio} dành cho hồ sơ diễn giả (UC41), chưa dùng ở nhóm này. Bảng
\code{UserToken} dùng chung cho ba mục đích (Bảng~\ref{tab:tk-token}); cột \code{TokenHash} chỉ chứa
bản băm, không bao giờ chứa chuỗi token gốc.

\begin{table}[H]
  \centering
  \caption{Các mục đích của token}
  \label{tab:tk-token}
  \begin{tabularx}{\textwidth}{lcLL}
    \toprule
    \textbf{Mục đích} & \textbf{Hiệu lực} & \textbf{Tạo khi} & \textbf{Hết hiệu lực khi} \\
    \midrule
    \code{VerifyEmail}   & 24 giờ  & Đăng ký, gửi lại email & Đã dùng; có token mới hơn \\
    \code{ResetPassword} & 30 phút & Quên mật khẩu & Đã dùng; có token mới hơn \\
    \code{Refresh}       & 7 ngày  & Đăng nhập, làm mới phiên & Làm mới; đăng xuất; đặt lại mật khẩu \\
    \bottomrule
  \end{tabularx}
\end{table}

\subsection{Thiết kế API}

Mọi API của nhóm nằm dưới đường dẫn \code{/api/auth} (Bảng~\ref{tab:tk-api}). Lỗi được trả theo chuẩn
ProblemDetails (RFC 9457) kèm trường \code{code} là mã lỗi nghiệp vụ (Bảng~\ref{tab:tk-ma-loi}); giao
diện dùng mã này để hiển thị thông báo đúng ngôn ngữ người dùng đang chọn.

\begin{table}[H]
  \centering
  \caption{Các API của nhóm Tài khoản}
  \label{tab:tk-api}
  \begin{tabularx}{\textwidth}{lp{4.2cm}cL}
    \toprule
    \textbf{Phương thức} & \textbf{Đường dẫn} & \textbf{Giới hạn} & \textbf{Kết quả thành công} \\
    \midrule
    POST & \code{/register}            & 5/phút & 201, email vừa đăng ký \\
    POST & \code{/verify-email}        & --     & 200, trạng thái mới của tài khoản \\
    POST & \code{/resend-verification} & 3/phút & 202 (luôn như nhau) \\
    POST & \code{/forgot-password}     & 3/phút & 202 (luôn như nhau) \\
    POST & \code{/reset-password}      & 5/phút & 204 \\
    POST & \code{/login}               & 5/phút & 200, access token + người dùng; cookie refresh \\
    POST & \code{/refresh}             & --     & 200, như đăng nhập; cookie mới \\
    POST & \code{/logout}              & --     & 204; xóa cookie \\
    GET  & \code{/me}                  & --     & 200, thông tin và quyền của người đang đăng nhập \\
    \bottomrule
  \end{tabularx}
\end{table}

\begin{table}[H]
  \centering
  \caption{Mã lỗi nghiệp vụ của nhóm Tài khoản}
  \label{tab:tk-ma-loi}
  \begin{tabularx}{\textwidth}{lcL}
    \toprule
    \textbf{Mã lỗi} & \textbf{HTTP} & \textbf{Khi nào} \\
    \midrule
    \code{Validation}         & 400 & Dữ liệu sai quy tắc (kèm lỗi theo từng ô) \\
    \code{EmailTaken}         & 409 & Email đã tồn tại \\
    \code{TokenInvalid}       & 422 & Token trong liên kết email không hợp lệ hoặc hết hạn \\
    \code{InvalidCredentials} & 401 & Sai email hoặc mật khẩu \\
    \code{AccountUnverified}  & 403 & Chưa xác thực email \\
    \code{AccountPending}     & 403 & Đang chờ quản trị viên duyệt \\
    \code{AccountLocked}      & 403 & Tài khoản bị khóa \\
    \code{SessionExpired}     & 401 & Refresh token không còn hiệu lực \\
    \code{Unauthorized}       & 401 & Thiếu hoặc sai access token \\
    \code{Forbidden}          & 403 & Đã đăng nhập nhưng thiếu quyền \\
    \code{TooManyRequests}    & 429 & Vượt giới hạn tần suất \\
    \bottomrule
  \end{tabularx}
\end{table}

Các chức năng gửi lại email và quên mật khẩu luôn trả 202 với cùng nội dung, kể cả khi email không
tồn tại hoặc đang trong 60 giây chờ; trả một mã khác cho trường hợp ``vui lòng chờ'' sẽ để lộ việc
email đó có tài khoản (BR08, BR09).
